\input{Laborbericht_header}%Datei mit Formatangaben, Paketen etc. für "Experten"
%Template für Laborberichte Werkstoffe III
%Author: Falk Wittel, ETH Zürich; fwittel@ethz.ch
%Herbstsemester 2019
%
%Für englische Berichte einfach deutsche Zeilen im Titelblatt auskommentieren und englische verwenden. Zudem ist in der Datei Laborbericht_header.tex die 4. Zeile auszukommentieren.

\begin{document}

% Anfang der Titelseite________________________________________________________________________________
\begin{titlepage}
	\flushleft
	{\scshape\Large Werkstoffe III \hspace{2.5cm} Laborbericht \hspace{2.5cm}HS 2022 \par}
	%{\scshape\Large Materials III \hspace{2.5cm} Lab report \hspace{2.5cm}HS 2019 \par}
	\vspace{6cm}
	{\huge\bfseries Titel des Labors\par}
	{\LARGE\bfseries Untertiel des Labors\par}
	\vspace{2cm}
    {\LARGE\bfseries Gruppe: X\par}
    {\LARGE\bfseries Datum des Labors: dd.mm.yyyy\par}
   % {\LARGE\bfseries Group: X\par}
   % {\LARGE\bfseries Date of lab: dd/mm/yyyy\par}
	\vspace{1.5cm}
	{\large Individuelle Beiträge zum Bericht in \%\par}
%	{\large Individual contributions to the report in \%\par}
\begin{table}[h]
\begin{tabular}{|l|l|l|l|l|l|l|l|}\hline
   & Student           & Konzep- & Aus- & Ver- & Korrek- & individueller & Unterschrift \\
   &             &  tion & wertung & -fassen & turlesen &  Anteil {[}\%{]} & \\ \hline
%   & Group           & concep- & data & editing & proof- & individual & signature \\
%   & member            &  tion & analysis &  & reading &  contribution {[}\%{]} & \\ \hline
   1  & Mustermann, Max   &            &                 &                   &                &                               &             \\ \hline
2  & Doe, Jane         &            &                 &                   &                &                               &             \\ \hline
3  & Musterfrau, Maria &            &                 &                   &                &                               &             \\ \hline
4  & Nordmann, Ola     &            &                 &                   &                &                               &             \\ \hline 
5  & N.N.              &            &                 &                   &                &                               &             \\ \hline
6  & N.N.              &            &                 &                   &                &                               &             \\ \hline
7  & N.N.              &            &                 &                   &                &                               &             \\ \hline
8  & N.N.              &            &                 &                   &                &                               &             \\ \hline
9  & N.N.              &            &                 &                   &                &                               &             \\ \hline
10 & N.N.              &            &                 &                   &                &                               &             \\ \hline
   &                   & 100\%      & 100\%           & 100\%             & 100\%          &                               &            \\ \hline
\end{tabular}
\end{table}
\vfill
% Bottom of the page
	{\large {\itshape Gruppenmitglieder deklarieren ihren Beitrag am Bericht und mit ihrer Unterschrift, das der Bericht selbständig und in eigenen Worten verfasst wurde und den Prinzipien wissenschaftlicher Redlichkeit folgt.} \Large{Unterzeichnen nicht vergessen!} \par}
%	{\large\itshape Group members declare their contributions to the report to be subject to all the rules governing academic honesty.  \bf{Do not forget to sign!} \par}
%Was unter Plagiat zu verstehen ist findet sich hier: https://ethz.ch/studierende/de/studium/leistungskontrollen/plagiate.html
\end{titlepage}
%%%%%%%%%%%%%%%Ende der Titelseite

\pagenumbering{Roman} % general chapters with Roman pagenumbers

%Abstract________________________________________________________________________
\chapter*{Zusammenfassung} \addcontentsline{toc}{chapter}{Zusammenfassung}
Kondensierte Version des gesamten Berichtes (ca. 200 Wörter). Gleiche Reihenfolge: Problem, Hypothese, Methoden, Kernergebnisse, Schlussfolgerung - je ein Satz.
 \clearpage

% Inhaltsverzeichnis_____________________________________________________________
 \setcounter{tocdepth}{2}
 \tableofcontents
 \clearpage

% Chapters______________________________________________________________________
\setlength{\parskip}{1em} 
\pagenumbering{arabic}				% Begin arabic page numbering (1,2,...)

% Einleitung_____________________________________________________________________
\chapter{Einleitung}\label{sec:}\markright{\thechapter~Einleitung}
%3 Absätze, Gesamtlänge ca. 1 Seite / 500 Wörter.
%Was ist das Problem?
Berichte schreiben will gelernt sein. Studierende haben oft Probleme mit der Erstellung von Berichten für Labore, Bachelorarbeiten und ähnliches. Diese Probleme sind nicht verwunderlich, denn es gibt in der Ingenieursausbildung leider kaum die notwendige Unterstützung und Lehre zum Erwerben der Schreibkompetenz. 

%Wieso ist es wichtig - was wurde gemacht?
Berichte sind essentiell um erworbenes Wissen zu dokumentieren und weiterzugeben. Eine  möglichst einheitliche Struktur und Stringenz eines Berichts sind wichtig, um eine Leserfreundlichkeit und ein effizientes Extrahieren des Inhalts zu gewährleisten. Innerhalb einer Fachdisziplin haben sich Konventionen zu Struktur und Format eines Berichts herausgebildet.

%Was wird von uns gemacht?
Im Rahmen des Werkstoffe III Labors erhalten Studierende die Gelegenheit mit Unterstützung durch Tutoren und Laborleiter das Erstellen von Berichten zu erlernen. Diese \LaTeX -Vorlage soll ebenfalls helfen. Sie bietet ein generisches, typisches Format eines Berichts, welches auch für die Bachelor- und Masterarbeit verwendet werden kann. In dieser Vorlage werden die wichtigsten \LaTeX-Elemente beispielhaft mit korrekter Formatierung gezeigt um den Einstieg in \LaTeX ~ mit Overleaf zu erleichtern. Es wird Hyerref verwendet, wodurch pdf-Dateien intern verlinkt sind.
% M&M______________________________________________________________________
\chapter{Material und Methoden}\label{chap:mm} \markright{\thechapter~Material \& Methoden}
Was wurde im Labor gemacht und wie wurde es gemacht. Wichtige Grundlagen müssen erklärt sein. Technische Aspekte des experimentellen Vorgehens, Vorrichtungen, Methoden der Datenerfassung und Analyse knapp dargestellt aber so vollständig, damit ein Leser die Versuche reproduzieren kann.
%%%%%%
\section{Kapitelhierarchie und Dokumentorgnisation}\label{sec:Struktur}
Es gilt: Chapter>Section>Subsection>Subsubsection>Paragraph, wobei bei einem kurzen Bericht selten auf die 3. Ebene, fast nie auf die 4. Ebene runter gegangen wird
%%%
\subsection{Subsection-Beispiel}\label{sec:Struktur-Subsection}
Das ist eine Subsection.
%%%%
\subsubsection{Subsubsection-Beispiel}\label{sec:Struktur-Subsection-Subsubsection}
Das Ist eine Subsubsection. Es ist die 4. Hierarchieebene, weshalb auf Darstellung der Nummer verzichtet wird. In einem kurzen Bericht sollte man nicht so tief gehen.
%%%%
\paragraph{Paragraph-Beispiel} kann referenziert werden und notfalls auch im Inhaltsverzeichnis aufgenommen werden (tocdepth=5). Es gibt keinen Zähler für Paragraphen.
%%%%%%%%%%%%
\section{Schriftgrössen definieren} \label{sec:Schriften}
Schriftgrössen sind. {\tiny Winzig} {\footnotesize Fussnote} {\small Kleiner} {\normalsize Normal} {\Large Grösser} {\huge riesig} {\Huge Gigantisch}. Soll Text betont werden, wird er über \textit{kursiv gesetzt}. Fetter Text ist \textbf{mit textbf.} Kommentiert wird mit \%. 

Schriften sollten im Header oder Stylefile geändert werden, da wir Inhalte und Formatierung in Latex trennen möchten. Grundsätzlich gibt es keine Vorgaben, nur sollten im Dokument verwendete Schriften durchgängig sein.
%%%%%%%%
\section{Formeln-Beispiele und Umgebungen}\label{sec:Formeln}
Formeln können im Text mit Dollarzeichen über $5=2>a\cdot \alpha$ ohne Zähler eingebettet werden. Formeln haben einen besonderen Zeichensatz, weshalb man Formelsymbole immer zwischen \$ \$ setzen sollte. Eine Formel kann natürlich auch abgesetzt werden wie die hier:
\begin{equation}\label{eq:Beispielformel}
P_\text{3.O} = a_0 +a_1\omega + a_2 \omega^2 +a_3\omega^3.
\end{equation}
Diese werden durchnummeriert (ausser equation*) und können über das Label referenziert werden und zwar als Gl.~\ref{eq:Beispielformel}. Da hyperref verwendet wird erzeugt dies im Pdf-Dokument einen Link. Achtung immer daran denken dass Formeln Teil des Satzes sind, also mit Punkt und Komma. Formeln sehen in Tex einfach super aus. Die Formelsyntax ist in Latex allerdings gewöhnungsbedürftig. Man muss entweder nachschauen (z.B. Ref.~\cite{wiki_book,BeuermannSascha2002EvlP}) oder kann auch andere Formeleditoren verwenden und dann in Latex konvertieren (z.B. Mathtype). Mehrzeilige Formeln benötigen die eqnarray Umgebung.
%%%%%%%%%%%%
\section{Tabellen-Beispiele und Umgebungen}\label{sec:Tabellen}
Tabellen sind in Latex etwas gewöhnungsbedürftig. Es gibt unterschiedliche Tabellenumgebungen (tabular, longtable, table, tabbing, etc.). Tabellen sind wie Abbildungen Floating-objekte wie auch Bilder, werden also vom Algorithmus im Dokument platziert. Man macht Spaltendefinitionen, die Beschriftung (caption) und dann den Inhalt. Spaltendefinitionen sind l=links; c=zentriert; r=rechts, p\{Spaltenbreite\}, $\|=$ senkrechte Tabellenlinie.  Ein Beispiel:%Verbatim  erlaubt es code zu drucken ohne dass es interpretiert wird.
\begin{verbatim} 
\begin{table}[H]
\centering
   \caption{Beispiel einer einfachen Tabelle.}   \label{tab:Beispiel}
\begin{tabular}{|p{60mm}|l|}
\hline
links oben               &     rechts oben  \\ \hline \hline
links unten   & rechts unten   \\ \hline
\end{tabular}
\end{table}
\end{verbatim}
\begin{table}[H]
\centering
   \caption{Beispiel einer einfachen Tabelle.}   \label{tab:Beispiel}
\begin{tabular}{|p{60mm}|l|}
\hline
links oben                 &     rechts oben  \\ \hline \hline
links unten   & rechts unten   \\ \hline
\end{tabular}
\end{table}
Es gibt auch kleine Hilfsprogramme (Tabelleneditoren), die die Erstellung richtig grosser Tabellen erheblich erleichtern können als kostenlose Programme, zur Konversion ganzer Excell-sheets oder auch als online converter (z.B. \url{https://www.tablesgenerator.com}), aus denen man dann nur noch die \LaTeX -Passage ins Dokument kopiert. Die Tabellenbeschriftung muss {\scshape\Large immer über} die Tabelle (s. Tab.~\ref{tab:Beispiel}) und bestehen wie bei Abbildungen aus einem vollständigen Satz.
%%%%%%%%%
\section{Abbildungen-Beispiele und Umgebungen}\label{sec:Abbildungen}
Abbildungen sind wie Tabellen Floating-objekte, die nach Formatvorgaben verteilt werden. Anweisungen sind h=hier; t=top Seite; b=Bottom Seite; p=Platzierung auf eigener Seite. Mit [H] kann Platzierung erzwungen werden, was aber nicht von Vorteil ist. Die Grössenangabe wird über width oder height gemacht. Es gibt auch scale, angle. Angaben sind in pt, mm, cm, oder Anteile der Zeilenbreite (textwidth) oder Seitenbreite. Bilder so klein wie möglich halten und wenn möglich PDF Vektorgrafiken verwenden. PDF, JPG und PNG Formate sind problemlos. Ein erstes Beispiel ist in Abb.~\ref{fig:Beispielbild1} gezeigt.
 \begin{figure}[htb]
	\begin{center}
	\includegraphics[width=0.85\textwidth]{IMAGES/writing_tutorial_fig3.jpg}
    \caption{\label{fig:Beispielbild1}Beispielbild unterschiedlicher Betonmischungen C1 und C2. (a) C1 nach der Mischung, (b) C2 ohne Verflüssiger und (c) C2 mit Verflüssiger.}
  	\end{center}
\end{figure}
%%%%%%%%%
\section{Referenzen-Beispiele}\label{sec:referenzen}
Alle referenzierten Objekte (Formeln, Tabellen, Kapitel, Abbildungen) haben eigene Zähler. Überall sollte ein label gesetzt sein. Sinnvoll sind label, die was aussagen z.B. ein Label für Gleichungen als label\{eq:Mitternachtsformel\}, für Abbildungen label\{fig:Beispielbild1\}, Tabellen \{tab:Materialdaten\} u.s.w.. Wenn man dann referenziert dann erfolgt das mit dem Befehlt ref{eq:Mitternachtformel}. Ich würde immer gleich jedes definierte Objekt mit einem Label versehen, das macht das Arbeiten leichter. Bisher haben wir Gleichung~\ref{eq:Beispielformel}, Tabelle~\ref{tab:Beispiel}, Abbildungen~\ref{fig:Beispielbild1}-\ref{fig:Beispielbild3} in Kapitel~\ref{chap:mm} definiert. Es ist das hyperref Paket im header gesetzt, d.h. im fertigen PDF Dokument sind alle Verweise, Kapitel in der Inhaltsangabe und Quellenangaben als Hyperlinks eingebettet (farbig markiert).
%%%%%%%%%%%
\section{Aufzählungen-Beispiele und Umgebungen}\label{sec:itemize}
Aufzählungen sind ideal um komplizierte, umfangreiche Sachverhalte zu strukturieren. Sie können als bullit items oder Aufzählungen sein: Ein Beispiel für Aufzälungen:
\begin{itemize}
\item Erstes Item
\item[$\Omega$] zweites Item. spezielles Item zwischen eckigen Klammern.
\end{itemize}
Aufzählungen können auch geschachtelt sein in allen möglichen Kombinationen:
\begin{enumerate}
  \item Erster Punkt
  \begin{enumerate}
     \item Erster Unterpunkt des ersten Punkten
     \item Zweiter Unterpunkt des ersten Punkten
  \end{enumerate}
  \item Zweiter Punkt
  \begin{itemize}
     \item Erstes Item
     \item[$\Omega$] zweites Item. spezielles Item zwischen Klammern. \[\].
  \end{itemize}
\end{enumerate}

Hier \glqq Weisheiten\grqq~ zum Bericht schreiben als Aufzählung:
\begin{itemize}
    \item Sei dir bewusst, dass die Erstellung eines Berichtes mehrere Stunden braucht, das ist bei der Anzahl Kreditpunkte der Lehrveranstaltung eingerechnet.
    \item Denk dran, Berichte sollen informativ sein und die Sache/das Problem steht im Vordergrund (kein Ich/Wir).
    \item Sei klar bei den Zielen des Berichtes.
    \item Verwende einfache, kurze Sätze, aber keine Emotionen oder Umgangssprache.
    \item Sei exakt und strukturiert beim Layout.
    \item KISS = keep it short \& simple; konzentriere dich auf das wesentliche und verwende auch nur wesentliches.
    \item Niemals Ergebnisse/Beobachtungen mit Interpretationen vermischen. Das sind unterschiedliche Kapitel.
    \item Sei knapp bei der Beschreibung bekannten und zitiere hier (formal korrekt). Investiere die Zeit lieber bei den Schlussfolgerungen.
    \item Sei deiner Arbeit gegenüber kritisch und versuche von anderen zu lernen, aber kopiere {\scshape\Large nie} von \grqq Freunden\grqq.
\end{itemize}
%%%%%%%%%%%%
\section{Quellenangaben}
Zitieren in \LaTeX ist sehr einfach. Es werden alles Quellen in einer *.bib Datei in Form von Bib-items gesammelt. Aus vielen Webseiten und anderen Programmen (wie Mendeley) lassen sich direkt bib-Einträge exportieren und kopieren. Ein Beispiel für eine Buch, eine Webseite und einen Artikel (mehr Beispiele in der {\em bibliography.bib} Datei):
\begin{verbatim}
@book{Weibull_McCool_2012,
title = "{Using the Weibull Distribution}",
author = "J.I.McCool",
year = "2012",
publisher = "Wiley"
}

@misc{wiki_book,
   author = "{Wikipedia contributors}",
   title = "Latex - Featured WikiBook",
   year = "2019",
   url = "https://en.wikibooks.org/wiki/LaTeX",
   note = "[Online; accessed 23-Octovber-2019]"
 }

@article{BeuermannSascha2002EvlP,
publisher = {Institut für Baumechanik und Numerische Mechanik, Uni Hannover; },
year = {2002},
title = {Erstellung von leistungsfähigen PDF-Dokumenten mit LATEX und den Paketen 
hyperref sowie thumbpdf},
edition = {Version 2.10},
author = {Beuermann, Sascha},
url        = {https://www.research-collection.ethz.ch/handle/20.500.11850/146361}
}
\end{verbatim}
Im Text zitiert wird über den Befehl {\em cite} die Quellen \cite{Weibull_McCool_2012,wiki_book,BeuermannSascha2002EvlP,Leigh-Phoenix-Beyerlein_2000}. Das Literaturverzeichnis wird automatisch erstellt. mit {\em citet} erhält man den Autor. Bsp: In \citet{BeuermannSascha2002EvlP} findet sich ....

% Resultate ___________________________________________________________________________________________________
\chapter{Resultate}\label{chap:Resultate} \markright{\thechapter~Resultate}
%%%%%%%%%
Was wurde herausgefunden? Hier werden alle ausgewerteten Ergebnisse, aber \textbf{KEINE} Interpretationen oder Diskussion, dargestellt - nur Fakten. Darstellung der ausgewerteten Daten (quantitativ mit statistischer Auswertung und Fehlerabschätzung wenn möglich) und nicht der Rohdaten. Grosse Tabellen oder Abbildungen mit Rohdaten gehören in den Anhang. Essentieller Bestandteil ist eine textliche, vergleichende Beschreibung der wichtigsten Ergebnisse aus den Abbildungen und Tabellen.  
%%%%%%%
\section{Erstes Ergebnis}\label{sec:Firstresult}
%%%%%%%%%
\section{Zweites Ergebnis}\label{sec:Secondresult}
Bei der Verwendung von Graphen wie in Abb.~\ref{fig:Beispielbild2} unbedingt aus Lesbarkeit der Achsenbeschriftungen achten. Verwendete Schriftgrösse sollte etwa der Grösse im Text entsprechen. Immer auf Einheiten achten und Legenden wenn erforderlich.
 \begin{figure}[H]
	\begin{center}
	\includegraphics[width=0.6\textwidth]{IMAGES/writing_tutorial_fig1.pdf}
    \caption{\label{fig:Beispielbild2} Spannungs-Dehnungs-Diagramm der Proben 1-5. Spannungswerte sind wahre Spannungen auf Grundlage der aktuellen Querschnittsfläche.}
  	\end{center}
\end{figure}

% Diskussion ___________________________________________________________________________________________________
\chapter{Diskussion}\label{sec:Disskussion} \markright{\thechapter~Diskussion}
Was bedeutet das alles? Hier erfolgt die Interpretation / Einordnung der Ergebnisse. Es wird erklärt was die Daten bedeuten und wie die Ergebnisse zu interpretieren sind. Man greift die Fragestellung der Einleitung wieder auf (ohne Wiederholung) und ordnet die Ergebnisse in das bisherige Wissen ein. Zwecks Lesefreundlichkeit können hierfür die relevanten Ergebnisse kurz aufgegriffen werden, eine ausführliche Wiederholung ist zu vermeiden. 

In der Diskussion werden die verschiedenen Ergebnisse auch verknüpft interpretiert, daher wird die Diskussion weniger in Unterkapitel unterteilt. Absätze können und sollen enthalten sein.

Als letzte(s) Unterkapitel der Diskussion finden sich Schlussfolgerungen/Fazit oder ein Ausblick/Empfehlungen, usw. wieder. Je nach Format können dies auch eigene Kapitel sein. 

 \begin{figure}[H]
	\begin{center}
	\includegraphics[width=0.8\textwidth]{IMAGES/writing_tutorial_fig2.png}
    \caption{\label{fig:Beispielbild3} Flussdiagramm des iterativen Arbeitsablaufs beim Abfassen von Berichten. (Quelle: \cite{1p_tutorial})}
  	\end{center}
\end{figure}
Schreiben ist ein iterativer Prozess, der mit der Konzeption anfängt, mit dem Aufbau, also dem anfertigen des Skeletts der Arbeit weitergeht und in Schritt 3 nach Abb.~ \ref{fig:Beispielbild3} mit der Ausformulierung weiter geht. An Ende steht das Korrekturlesen. Zudem sollten sich die Autoren fragen, ob sie die Checkliste in Abb.~\ref{fig:Beispielbild4} wirklich erfüllen.
 \begin{figure}[H]
	\begin{center}
	\includegraphics[width=0.9\textwidth]{IMAGES/writing_tutorial_fig4.pdf}
    \caption{\label{fig:Beispielbild4} Checkliste und Grundlage des Bewertungsschemas für Berichte. (Quelle: \cite{1p_tutorial})}
  	\end{center}
\end{figure}

% Bibliography__________________________________________________________________
% Literature (Additional references can be added to the .bib-file manually, or by using, for example, the free application JabRef). 
\bibliography{bibliography}%Welche Quellen wurden verwendet? Ein Zitierknigge findet sich hier: https://ethz.ch/content/dam/ethz/main/education/rechtliches-abschluesse/leistungskontrollen/plagiat-zitierknigge.pdf

% Anhang______________________________________________________________________
\appendix
\input{Laborbericht_Anhang} %für umfangreiche Rohdaten oder Prozedurenbeschreibungen

\end{document}
